\documentclass[acmsmall,screen,review]{acmart}

\usepackage{proof}
\usepackage{fancyvrb}

\usepackage{latexsym}
\usepackage{tikz}
\usepackage{graphicx}
\usepackage{mathpartir}
\usepackage{color}
\usepackage{pgflibraryarrows}
\usepackage{pgflibraryshapes}
\usepackage{pgfbaseimage}
\usepackage{pgfplots}

\usetikzlibrary{decorations.text}
\usetikzlibrary{trees}
\usetikzlibrary{fadings}

\usepackage{xcolor}

\usepackage{amsthm}

\let\Bbbk\relax
\usepackage{amssymb}

\let\Bbbk\relax
\usepackage{amssymb}

\usepackage[ruled,vlined]{algorithm2e}
\SetAlgoLined
\SetAlgoNoEnd
\SetArgSty{textup}
\SetKwProg{Program}{Program}{\hspace{0.3em}:}{}
\SetKwInOut{Input}{input}\SetKwInOut{Output}{output}
\SetKwIF{Match}{ElseMatch}{Otherwise}{match}{with}{else}{else}{endmatch}

\usepackage{fancyhdr}

\usepackage{booktabs}

\usepackage{multicol}

\setcopyright{acmlicensed}
\copyrightyear{2026}
\acmYear{2026}
\acmDOI{XXXXXXX.XXXXXXX}

\acmJournal{JACM}
\acmVolume{37}
\acmNumber{4}
\acmArticle{111}
\acmMonth{8}

%%
%% Submission ID.
%% Use this when submitting an article to a sponsored event. You'll
%% receive a unique submission ID from the organizers
%% of the event, and this ID should be used as the parameter to this command.
%%\acmSubmissionID{123-A56-BU3}


%%
%% The majority of ACM publications use numbered citations and
%% references.  The command \citestyle{authoryear} switches to the
%% "author year" style.
%%
%% If you are preparing content for an event
%% sponsored by ACM SIGGRAPH, you must use the "author year" style of
%% citations and references.
%% Uncommenting
%% the next command will enable that style.
%%\citestyle{acmauthoryear}


\newcommand{\abs}[4]{{#1}\, #2\! : \! #3.\, #4}
\newcommand{\absu}[3]{{#1}\, #2.\, #3}
\mathchardef\mhyph="2D % Define a "math hyphen"

% unicode
\usepackage[utf8]{inputenc}
\DeclareUnicodeCharacter{3BB}{\ensuremath{\lambda}}
\DeclareUnicodeCharacter{3B4}{\ensuremath{\delta}}
\DeclareUnicodeCharacter{3C9}{\ensuremath{\omega}}
\DeclareUnicodeCharacter{3B3}{\ensuremath{\gamma}}
\DeclareUnicodeCharacter{3C4}{\ensuremath{\tau}}
\DeclareUnicodeCharacter{21D2}{\ensuremath{\To}}
\DeclareUnicodeCharacter{22B2}{\ensuremath{\lhd}}
\DeclareUnicodeCharacter{2080}{\ensuremath{_0}}
%\definecolor{darkred}{RGB}{100,10,10}


\newcommand{\BD}[1]{\textcolor{olive}{[BD: #1]}}
\newcommand{\PA}[1]{\textcolor{cyan}{[PA: #1]}}

\newcommand{\utp}[0]{\mathcal{U}}
\newcommand{\lin}[0]{\lambda^{\iota\nu}}
\newcommand{\interp}[1]{\llbracket #1 \rrbracket}
\newcommand{\interpl}[0]{\llbracket}
\newcommand{\interpr}[0]{\rrbracket}
\newcommand{\choice}[0]{\xi}
\newcommand{\simto}[0]{\triangleright}
\newcommand{\elcap}[0]{\cap}

\newcommand{\dedu}[0]{\triangleright}
\newcommand{\ctxtup}[1]{\ulcorner #1 \urcorner}
\newcommand{\ctxtdn}[1]{\llcorner #1 \lrcorner}
\newcommand{\ctxtmv}[1]{\langle #1 \rangle}

\newcommand{\case}[1]{\noindent\underline{Case #1:}}
\newcommand{\caseb}[0]{\noindent\underline{Case:}}

\newcommand{\pcase}[1]{\noindent\underline{#1:}}


\newcommand{\adbmal}[0]{\reflectbox{$\lambda$}}

\newcommand{\casetri}[0]{\RHD}
\newcommand{\lke}[0]{\lhd}

\newcommand{\lam}[2]{\lambda\, #1.\, #2}
\newcommand{\mal}[2]{\abdmal\, #1.\, #2}
\newcommand{\all}[2]{\forall\, #1.\, #2}
\newcommand{\Lam}[2]{\Lambda\, #1.\, #2}
\newcommand{\thereis}[2]{\exists\, #1.\, #2}
\newcommand{\nuu}[2]{\nu\, #1.\,#2}
\newcommand{\inn}[2]{\langle #1 \rangle_{#2}}
\newcommand{\exel}[2]{\adbmal\,#1.\,#2}
\newcommand{\recc}[2]{\textit{rec}\, #1.\, #2}
\newcommand{\mew}[2]{\mu_{#1}\, #2}
\newcommand{\convl}[0]{\blacktriangleleft}
\newcommand{\convr}[0]{\blacktriangleright}
\newcommand{\conjin}[0]{\smalltriangleleft}
\newcommand{\conjout}[0]{\smalltriangleright}
\newcommand{\dcdot}[0]{\mathbin{\ast}}
\newcommand{\To}[0]{\Rightarrow}
\newcommand{\releq}[0]{\topdoteq}

\newcommand{\omga}[4]{\omega\ #1(#2)\,:\,#3.\ #4}
\newcommand{\clause}[2]{#1 \to #2}

\newcommand{\hysep}[0]{\mathbin{\blacklozenge}}

\newcommand{\todo}[1]{\textbf{[#1]}}
\newcommand{\rref}[1]{\{\ref{#1}\}}

\newtheorem{theorem}{Theorem}
\newtheorem{lemma}[theorem]{Lemma}
\newtheorem{corollary}[theorem]{Corollary}
\newtheorem{definition}[theorem]{Definition}
\newtheorem{notation}[theorem]{Notation}
\newtheorem{proposition}[theorem]{Proposition}


%\newtheorem{definition}[theorem]{Definition}


% \newenvironment{mycomment}
%     {\color{blue} \it
%     }
%     {
%     }
\newcommand{\shift}[2]{{\uparrow^{#1}_{#2}}}
\newcommand{\textblue}[1]{\textcolor{blue}{#1}}

\newcommand{\tl}[0]{\mathsf{L}}
\newcommand{\tn}[0]{\mathsf{N}}
\newcommand{\subt}[4]{(#1)\ #2 <:^{#3} #4}

\newcommand{\st}{\mathrel{<:}}
\newcommand{\coe}{\mathsf{Coe}}
\newcommand{\ncoe}{\mathsf{NCoe}}
\newcommand{\nf}{\mathsf{nf}\,}
\newcommand{\F}{\mathsf{F}\,}
\newcommand{\id}{\mathsf{id}}

\newcommand{\cid}[1]{\mathsf{id}\ #1}
\newcommand{\ccomp}[2]{#1 \cdot #2}
\newcommand{\carr}{\mathsf{arr}}
\newcommand{\cmono}{\mathsf{mono}}
\newcommand{\cfmap}{\mathsf{fmap}}

\newcommand{\nid}{\mathsf{nid}}
\newcommand{\narr}{\mathsf{narr}}
\newcommand{\nmono}{\mathsf{nmono}}
\newcommand{\nfmap}{\cdot \mathsf{nfmap}}
\newcommand{\naarr}{\cdot \mathsf{nfmap}}

\newcommand{\tbool}{\mathsf{bool}}


%%
%% end of the preamble, start of the body of the document source.
\begin{document}

\title{Inserting Fmap with Coercive Subtyping}

%% \author{Ryan Doenges}
%% \email{me@ryandoeng.es}
%% \affiliation{%
%%   \institution{Boston College}
%%   \city{Boston}
%%   \state{Massachusetts}
%%   \country{USA}
%% }

%% \author{Caden Parajuli}
%% \email{parajulc@bc.edu}
%% \affiliation{%
%%   \institution{Boston College}
%%   \city{Boston}
%%   \state{Massachusetts}
%%   \country{USA}
%% }

%% \author{Aaron Stump}
%% \email{stumpaa@bc.edu}
%% \affiliation{%
%%   \institution{Boston College}
%%   \city{Boston}
%%   \state{Massachusetts}
%%   \country{USA}
%% }


%%
%% By default, the full list of authors will be used in the page
%% headers. Often, this list is too long, and will overlap
%% other information printed in the page headers. This command allows
%% the author to define a more concise list
%% of authors' names for this purpose.
\renewcommand{\shortauthors}{\ }%Blanchette et al.}

%%
%% The abstract is a short summary of the work to be presented in the
%% article.
\begin{abstract}
  This paper shows how to use subtyping to infer uses of the basic
  monad operations.
\end{abstract}

%%
%% The code below is generated by the tool at http://dl.acm.org/ccs.cfm.
%% Please copy and paste the code instead of the example below.
%%
\begin{CCSXML}
<ccs2012>
<concept>
<concept_id>10011007.10011006.10011008.10011009.10011012</concept_id>
<concept_desc>Software and its engineering~Functional languages</concept_desc>
<concept_significance>500</concept_significance>
</concept>
<concept>
<concept_id>10011007.10011006.10011008.10011024.10011033</concept_id>
<concept_desc>Software and its engineering~Recursion</concept_desc>
<concept_significance>500</concept_significance>
</concept>
</ccs2012>
\end{CCSXML}

\ccsdesc[500]{Software and its engineering~Functional languages}

\keywords{subtyping, functors}

%\received{20 February 2007}
%\received[revised]{12 March 2009}
%\received[accepted]{5 June 2009}

\theoremstyle{definition}
\newtheorem{specification}[section]{Specification}

%%
%% This command processes the author and affiliation and title
%% information and builds the first part of the formatted document.
\maketitle

\section{Introduction}
Typeclass abstractions like Functor, Applicative, and Monad are essential to
modern functional programming.
%
The liberal application of type class combinators allows
concise programming even in the presence of 
complex effects and data structures.
%
For example, consider writing a Haskell program that increments integers wrapped
in a \verb|Maybe| type.
\begin{Verbatim}
incrMaybe : Maybe Int -> Maybe Int
incrMaybe (Just x) = Just (succ x)
incrMaybe Nothing  = Nothing
\end{Verbatim}
This naive definition of \verb|incrMaybe| has to unwrap its argument with
a pattern match and then wrap its output back up.
A more idiomatic definition uses \verb|<$>|, the fmap combinator, which
handles the unwrapping and wrapping.
\begin{Verbatim}
incrMaybe x = succ <$> x
\end{Verbatim}
In the presence of sufficient type information, we are interested in further
reducing the boilerplate involved by making fmap implicit. This permits writing
\verb|incrMaybe| with no use of fmap at all.
\begin{Verbatim}
incrMaybe x = succ x
\end{Verbatim}
It seems as if there is only one canonical way to make this program
make sense, and automatically lifting functions to operate on arrays and other
data structures is already prevalent both in programming languages and
mathematical notation.

In this paper, we study a \emph{coercive subtyping} discipline
that infers where to insert implicit fmaps using subtyping.
Each derivation of \(A \st B\) then corresponds to a coercion \(A \to B\)
expressed in a simple combinator calculus.
In coercive subtyping, type inference is coupled to an elaboration pass,
which inserts coercions in the program based on the results of subtyping
queries.
This makes the operational meaning of surface programs depend on the result of
type inference, which breaks with a more naive view of inference as
merely tacking missing type annotations on to a program with already fixed
semantics.
Indeed, many programming languages confound their users with strange behavior
involving implicit coercions, whether they are inserted statically or
dynamically.

We can avoid these confusions by establishing \emph{coherence}, a basic
metatheoretic property of systems of coercions. This problem was studied as
early as 1980 by Reynolds~\cite{reynolds80}, who outlined an algebraic approach
to proving the coherence of coercions. In his view, coercions form
a \emph{category}: there are types \(A, B, \dots\) and between any two types
there is a set of coercions \(\coe(A, B)\). There are always identity coercions
\(\id_A \in \coe(A, A)\), which do nothing, and there are associative and unital
composition operations \(\circ : \coe(B, C) \times \coe(A, B) \to \coe(A, C)\).
The coercions are expected to form a preorder, that is, for every pair of arrows
\(c1, c2 \in \coe(A, B)\), \(c1 \equiv c2\). He also studies their interaction
with generic operators of the language, but we will restrict our attention to
the coercions themselves.

As is typical in programming languages, we establish our coercions are coherent
by showing every coercion \(c : A \subty B\) can be put into a normal form \(\nf
c\) that is equivalent to \(c\), and then show that normal forms are unique.
Rather than relying on rewriting, we adopt the technique of \emph{normalization
by evaluation} (NbE).

Besides giving some confidence in the safety of this programming model, the
proof of soundness has independent interest as a case study in the application
of normalization by evaluation to implicit coercions.

\section{Overview}
The types and subtyping rules at issue are shown in Figure~\ref{fig:calculus}.
We call these the declarative rules (as opposed to the algorithmic rules).
The only nonstandard rule is \textsc{Fmap}, which permits using functions \(A
\to B\) at the lifted type \(\F A \to \F B\). This rule may be iterated,
permitting the use of functions at \(\F A \to \F B\), \(\F \F A \to \F \F B\),
and so on. These rules serve as a good specification of the system, but they
have practical and theoretical limitations.

The declarative rules are not suitable for an implementation because they are not
deterministic: too many rules apply to possible subtyping queries.
Besides the fact that it is always possible to appeal to \textsc{Trans},
queries of the form \(A \to B \subty \F U \to \F V\)  must choose between
\textsc{Arr} and \textsc{Fmap}.

They are also questionable from a semantic point of view because a single
judgment may have many derivations and the coercions associated with these
derivations may not be equivalent.

\begin{figure}
  \begin{minipage}{0.25\textwidth}
  \[
  \begin{array}{rcl}
    A & ::= & \tbool \\
    & | & A \to A \\
    & | & \F A
  \end{array}
  \]
  \end{minipage}%
  \begin{minipage}{0.75\textwidth}
  \begin{mathpar}
    \inferrule*[Left=Refl]{ }{A \subty A}
    
    \inferrule*[Left=Trans]{A \subty B \\ B \subty C}{A \subty C}
   
    \inferrule*[Left=Arr]{U \subty A \\ B \subty V}{A \to B \subty U \to V}
 \\ 
    \inferrule*[Left=Mono]{A \subty B}{\F A \subty \F B}

    \inferrule*[Left=Fmap]{ }{A \to B \subty \F A \to \F B}
  \end{mathpar}
  \end{minipage}
  \caption{Types and subtype relationships for a calculus with implicit fmap.}\label{fig:calculus}
\end{figure}

\todo{Describe the design of normalized coercions by constructing
``multiplication tables''}

\section{Coercions}
To define coercions, we introduce proof terms to the subtyping judgment and then
define an equivalence on proof terms. In the new judgment \(c : A \subty B\),
the proof term \(c\) is an intrinsically typed coercion; for proof rules that
consume multiple derivations the coercion \(c\) is defined using a coercion
combinator (Fig.~\ref{fig:coe}).

\begin{definition}[Coercion congruence]
  The equivalence relation \(\equiv\) is the least congruence \(\equiv\) such
  that for all \(c : A \subty B\),
\begin{align*}
\ccomp{\cid{}}{c} &\equiv c \\
\ccomp{c}{\cid{}} &\equiv c \\
\ccomp{c_1}{(\ccomp{c_2}{c_3})} &\equiv \ccomp{(\ccomp{c_1}{c_2})}{c_3} \\
\cmono\ (\cid{A}) &\equiv \cid{\F A} \\
\cmono\ (\ccomp{c_1}{c_2}) &\equiv \ccomp{\cmono\ c_1}{\cmono\ c_2} \\
\carr\ \cid{A}\ \cid{B} &\equiv \cid{A \to B} \\
\carr\ (\ccomp{c_1}{c_2})\ (\ccomp{c_3}{c_4}) &\equiv
\ccomp{\carr\ c_2\ c_3}{\carr\ c_1\ c_4} \\
\ccomp{\cfmap}{\carr\ (\cmono\ c_1)\ (\cmono\ c_2)} &\equiv \ccomp{\carr\ c_1\ c_2}{\cfmap}
\end{align*}
\end{definition}
In order, these laws say that coercions form a category, \(\cmono\) is a functor
on that category, and \(\carr\) is a bifunctor on that category.  The final law
about \(\cfmap\) is more complicated becuase \(\cfmap\) is not a functor on
coercions but instead a functor on types. In the equational theory of coercions,
all that can be said about \(\cfmap{}\) is how it interacts with other coercions on
arrow types.

\subsection{Compiling coercions correctly}
In order to justify the equational theory and show that these coercion
combinators can be soundly compiled to a programming language, we provide a
translation into typed lambda calculus equipped with an abstract functor
\texttt{F}.  Then we prove that the translation preserves equivalence, which
is given by the equational theory in the source language and by
\(\beta\eta\)-equivalence and fmap laws in the target.
\begin{minipage}{0.5\textwidth}
\begin{align*}
  \cc\tbool &= \ltbool \\
  \cc{A \to B} &= \cc{A} \ltarr \cc{B} \\
  \cc{\F A} &= \ltf \cc{A} \\
\end{align*}
\end{minipage}
\begin{minipage}{0.5\textwidth}
\begin{align*}
  \cc{\cid{A}} &= \lambda x : \cc{A}.\, x \\
  \cc{\ccomp{c_1}{c_2}} &= \lambda x.\,\cc{c_2}\ (c_1\ x) \\
  \cc{\carr\ c_1\ c_2} &= \lambda f\, x.\,c_2\ (f\ (c_1\ x)) \\
  \cc{\cmono\ c} &= \ltfmap\ \cc{c} \\
  \cc{\cfmap} &= \ltfmap
\end{align*}
\end{minipage}

\begin{proposition}[\ccname{} sound \agdachk{interp-sound}]
  For all \(c_1, c_2 : A \subty B\), \[
    c_1 \equiv c_2 \Longrightarrow \cc{c_1} \equiv_{\beta\eta\ltf} \cc{c_2}
  \]
\end{proposition}

\subsection{Normal forms for coercions}

\begin{figure}
  \begin{flushleft}
    \fbox{\(c : A \subty B\)}
  \end{flushleft}
  \begin{mathpar}
    \inferrule*[Left=CId]{ }{\cid{A} : A \nsubty A}
    
    \inferrule*[Left=CComp]
      {c_1 : A \subty B \\ c_2 : B \subty C}
      {\ccomp{c_1}{c_2} : A \subty C}
  \\ 
    \inferrule*[Left=CArr]
      {c_1 : U \subty A \\ c_2 : B \subty V}{\carr\ c_1\ c_2 : A \to B \subty U \to V}
 \\ 
    \inferrule*[Left=CMono]{c : A \subty B}{\cmono\ c : \F A \subty \F B}

    \inferrule*[Left=CFmap]{ }{\cfmap : A \to B \subty \F A \to \F B}
  \end{mathpar}
  \caption{Declarative coercions.}\label{fig:coe}
\end{figure}

\begin{figure}
  \begin{flushleft}
    \fbox{\(n : A \nsubty B\)}\quad\fbox{\(\mathit{na} : A,B  \nasubty U, V\)}
  \end{flushleft}
  \begin{mathpar}
    \\
    \inferrule*[Left=NId]{ }{\nid : \tbool \nsubty \tbool}

    \inferrule*[Left=NMono]{n : A \nsubty B}{\nmono\ n : \F A \nsubty \F B}
    \\
    \inferrule*[Left=NArr]
      {\mathit{na} : A, B \nasubty U, V}{\narr\ \mathit{na} : A \to B \nsubty \F A \to \F B}

    \inferrule*[Left=NAFmap]
      {\mathit{na} : A, B \nasubty U, V}{\nfmap{\mathit{na}} : A, B \nasubty \F U, \F V}
      \\

    \inferrule*[Left=NAArr]
      {n_1 : U \nsubty A \\ n_2 : B \nsubty V}{\naarr\ n_1\ n_2 : A, B \nasubty U, V}
  \end{mathpar}
  \caption{Normalized coercions.}\label{fig:calculus}
\end{figure}

We would like to prove the uniqueness of normal coercions.  This ought to look like an induction on one
paired with some inversion lemmas on the other, but there are issues
in the arrow case. For instance, inverting \(\na_2 : A, B \nasubty \F U, \F V\)
suggests that \(na_2\) has the form \(\na_2 = \nfmap{\na'}\) or
\(\na_2 = \naarr\ n_1'\ n_2'\), but both possiblities remain plausible.
Fortunately, there are some invariants of subtyping that will allow ruling out
the latter possibility.

The subtyping relations \(A \nsubty B\) and \(A, B \nasubty U, V\) exhibit
predictable relationship between the number of leading functors in
the types on their left and right. To this point, define the \emph{peel number} of a
type to be  \(\peel{\F A} = 1 + \peel A\) and \(\peel A = 0\) otherwise.  The
peel number is preserved by normal subtyping and increased by arrow subtyping.
\begin{proposition}[Peel invariants]
  Let \(A, B, U\), and \(V\) be types.
  \begin{enumerate}
    \item If \(A \nsubty B\) then \(\peel A = \peel B\).
    \item If \(A, B \nsubty U, V\) then \(\peel A + k = \peel U\) and \(\peel B + k = \peel V\) for some \(k\).
  \end{enumerate}
\end{proposition}
\begin{proof}
  By mutual induction on \(\nsubty\) and \(\nasubty\).
\end{proof}
With this in hand, we can distinguish arrow coercions. In fact, the \(k\) that
appears in the invariant for arrow subtyping counts the number of
\(\mathsf{nfmap}\) constructors that appear in the corresponding derivation.

\begin{proposition}[Uniqueness of normal forms \agdachk{ncoe-uniq}]\label{thm:uniqueness}
  Let \(A, B, U\), and \(V\) be types.
  \begin{itemize}
    \item If \(n : A \nsubty B\) and \(n' : A \nsubty B\) then \(n = n'\).
    \item If 
      \(\na : A, B \nasubty U, V\)
      and
      \(\na' : A, B \nasubty U, V\)
      then 
      \(\na = \na'\).
  \end{itemize}
\end{proposition}
\begin{proof}
  By mutual induction on \(n\) and \(\na\). We show a few representative cases.
\begin{itemize}
  \item \todo{mono case}
  \item \todo{na-fmap case}
\end{itemize}
\end{proof}

\section{Coherence}
In this section, we prove that all coercions of the same type are equivalent. We
establish our result using the methods of normalization by evaluation, which is
more commonly used to establish that a language has decidable equality by
lifting it from a model with an already decidable relation, which is often
closer to syntactic equality. In this section we establish coherence for
coercions by lifting it from a proof of uniqueness for normal coercions.

\begin{theorem}
  If \(c_1, c_2 : A \subty B\) then \(c_1 \equiv c_2\).
\end{theorem}

\begin{itemize}
  \item The equational theory of coercions: congruence closure of category laws
    plus functor/bifunctor laws and the ``special fmap law"
    \begin{itemize}
      \item informal justification in terms of fmap laws and associativity of
        composition (for terminating programs.... if we are going to
        be careful...)
      \item explanation of why this is necessary at all
    \end{itemize}
  \item The equational theory of normalized coercions: just syntactic equality.
\end{itemize}

\subsection{Reification and evaluation}
Normal coercions are immediately interpretable as coercions, as exhibited by the
following straightforward definition of \(\reify\).
\begin{align*}
  \reify\ \nid &= \cid{\tbool} \\
  \reify\ (\nmono\ n) &= \cmono\ (\reify\ n) \\
  \reify\ (\narr\ \na) &= \reifyna \na \\
  \reifyna\ (\naarr\ n_1\ n_2) &= \carr\ (\reify\ n_1)\ (\reify\ n_2) \\
  \reifyna\ (\nfmap{\na}) &= \ccomp{\reifyna \na}{\cfmap}
\end{align*}
Note that the reification of \(\nfmap{\na}\) puts \(\cfmap\) after the
reification of \(\na\) and not before. This is an arbitrary choice, like the
choosing between cons lists and snoc lists. Putting \(\cfmap\) at the front
would affect the ``bias'' of subsequent proofs but all the top-level theorems
would say the same things.

Reification is sound because all functions preserve equality. By composition,
\(\interp{-} \circ \reify\) is also sound. This means our syntactic equality
agrees with the intended model in lambda calculus.  The type of reify is also
informative. Read through Curry-Howard, it says that the normal subtyping
relation is sound (but not necessarily complete) for subtyping. \[
  A \nsubty B \Longrightarrow A \subty B
\]

Evaluation, which passes from coercions to normal coercions, relies on some
definable operators on normal coercions. This is where most of the work happens.
\[
  \mkid{A} : A \nsubty A \qquad \mkcomp : A \nsubty B \times B \nsubty C \to A \nsubty C
\]
These correspond to standard lemmas in the metatheory of subtyping.
The first proves a reflexivity rule is admissible,
while the second proves a transitivity rule is admissible.
Evaluation uses these to establish the completeness of normal subtyping.
Promoting these results from theorems to functions makes them proof-relevant,
which will ultimately allow us to show a stronger result: implicit coercions in
both theories behave the same.

The identity and composition operators permit a structurally recursive
definition of \(\eval\), which is otherwise defined in terms of constructors.
From an algebraic point of view, we are exhibiting an algebraic structure on
normal coercions that then induces the interpretation \(\eval\) from the free
algebra---coercions---into normal coercions.
\begin{align*}
  \eval{\cid{A}} &= \mkid{A} \\
  \eval{(\ccomp{c_1}{c_2})} &= \mkcomp\ (\eval{c_1})\ (\eval{c_2}) \\
  \eval{(\cmono\ c)} &= \nmono\ (\eval{c}) \\
  \eval{(\carr\ c_1\ c_2)} &= \narr\ (\naarr\ (\eval{c_1})\ (\eval{c_2})) \\
  \eval{\cfmap} &= \narr\ \nfmap \\
\end{align*}

The identity operator is defined by induction on type. This helps explain why
the normal coercions do not have a built-in \(\nid_A\) operator for any type
\(A\), but only include it for base types. Otherwise, at every type \(A\) the
two distinct proofs \(\nid_A\) and \(\mkid{A}\) would both prove
\(A \nsubty A\), which would break uniqueness (Prop.~\ref{thm:uniqueness}).
\begin{align*}
  \mkid{\tbool} &= \nid \\
  \mkid{A \to B} &= \narr\ (\naarr\ (\mkid{A})\ (\mkid{B})) \\
  \mkid{\F A} &= \nmono\ (\mkid{A})
\end{align*}

The composition operator is more complex. In order to design the composition
operator and, in fact, design the algebra of normal coercions,
\todo{multiplication table section}
\todo{show definition of mkcomp}

It is easy to show that evaluation preserves coercion equivalence since
evaluation preserves type and normal coercions are unique.
This is remarkable because it implies (for instance) that our defined
composition and reflexivity operators obey all the laws of coercions.  Given the
complexity of the composition operator, this is a very simple proof.
\begin{proposition}[Soundness of evaluation \agdachk{eval-sound}]
  For all \(c_1, c_2 : A \subty B\),
  if \(c_1 \equiv c_2\) then \(\eval{c_1} = \eval{c_2}\).
\end{proposition}
\begin{proof}
  Since \(\eval{c_1}, \eval{c_2} : A \nsubty B\), this follows from
  uniqueness (Prop.~\ref{thm:uniqueness}).
\end{proof}

\subsection{Inverting evaluation}
We now hew close to the path laid by normalization by evaluation and prove that
\(\reify\) is a left inverse to \(\mathsf{eval}\). This makes the normalization
function \(\reify \circ \mathsf{eval}\) sound. The proof of this lemma exercises
all the rules of the equivalence relation \(\equiv\) on coercions,
including the ``special fmap rule.''

The first order of business in this proof is to establish that \(\reify\) is a
homomorphism for the defined algebra on normal coercions.
\begin{proposition}[\agdachk{reify-mkid}]\label{reify-mkid}
  For all \(A\), \(\reify\ \mkid{A} \coeeq \cid{A}\).
\end{proposition}
\begin{proof}
  By induction on \(A\) using the identity laws for \(\cmono\) and \(\carr\).
\end{proof}

The lemma about composition is really two lemmas, since the composition operator
also has to compose arrow coercions.
\begin{proposition}[\agdachk{reify-mkcomp}]\label{reify-mkcomp}
  Let \(A, B, U, V, X\), and \(Y\) be types.
  \begin{enumerate}
    \item For all \(n_1 : A \nsubty B\) and \(n_2 : B \nsubty C\),
  \[\reify\ (\mkcomp\ n_1\ n_2) \equiv \ccomp{(\reify\ n_1)}{(\reify\ n_2)}.\]
    \item For all \(\na_1 : A, B \nasubty U, V\) and \(\na_2 : U, V \nasubty X, Y\), \dots \todo{mkcompna}
    \item For all \(\na : A, B \nasubty U, V\), \dots \todo{na-fmap-l}
  \end{enumerate}
\end{proposition}

Now we can prove that reify inverts evaluation.
\begin{lemma}[Reify-eval \agdachk{reify-eval}]\label{reify-eval}
  If \(c : A \subty B\), then \(\reify\ \eval{c} \equiv c\).
\end{lemma}
\begin{proof}
  By induction on \(c\).
  \begin{itemize}
    \item Case \(c = \cid{A}\): By Prop.~\ref{reify-mkid}.
    \item Case \(c = \ccomp{c_1}{c_2}\): By Prop.~\ref{reify-mkcomp} and the IH.
    \item Case \(c = \cmono\ c\): By the IH.
    \item Case \(c = \carr\ c_1\ c_2\): By the IH.
    \item Case \(c = \cfmap\): We reason equationally.
          \begin{align*}
            \cfmap \equiv ?
          \end{align*}
          \todo{finish proof}
  \end{itemize}
\end{proof}

\begin{theorem}
  If \(c_1, c_2 : A \subty B\) then \(c_1 \coeeq c_2\).
\end{theorem}
\begin{proof}
  We reason equationally since \(\equiv\) is an equivalence relation.  The
  middle step uses equality, which is finer than \(\equiv\) and respected by
  \(\reify\).
  \begin{align*}
    c_1 &\equiv \reify (\eval{c_1}) & \text{by symmetry of \(\equiv\) and normalization (Lemma~\ref{reify-eval})}&\\
    &= \reify (\eval{c_2}) & \text{by uniqueness (Prop~\ref{thm:uniqueness})} &\\
    &\equiv c_2 & \text{by normalization (Lemma~\ref{reify-eval})} &\qedhere
  \end{align*}
\end{proof}



\section{Implementation in Agda}
All of our definitions and theorems are proved in Agda,
where we made heavy use of dependently typed syntax and dependent pattern matching
to keep our theorems concise and our proofs simple.
The presentation of the system in the paper glosses over details that were
checked by Agda, like the well-typedness of the output of the reification and
evaluation functions.

\todo{key defns here in a syntax-highlighted environment}

We found Agda well suited to the task of mechanizing normalization by evaluation
for coercions. All proof assistants are adequate for algebra over a single sort,
but in many-sorted and indexed algebras good pattern matching becomes indispensable.
Agda's dependent pattern matching allowed us to use an intrinsically typed syntax
and write concise proofs. Recall the multiplication table from Section~\todo{ref
to mult discussion once it's written}. In Agda, many of the cases we explained
away are automatically ruled out by the pattern matching algorithm. In inductive
proofs, this makes a big difference in the usability of the system and the
efficiency of formalizing metatheory.

While reasoning about coercions, we rely on a user-defined equivalence relation
instead of Agda's propositional equality. While Agda has special ``rewrite
patterns'' for rewriting with equations on the left side of a pattern match, we
did not find the same level of built-in support for rewriting with custom
equivalences.  This means that our proofs end up containing big proof terms that
are littered with congruence and transitivity rules, like the following clause
from the proof of \texttt{reify-eval}. What this term really does is a series of
rewrites along \texttt{c-arr-fmap}, \texttt{reify-mkid}, \texttt{c-mono-id},
\texttt{c-arr-id}, and \texttt{c-id-r}, but it also has to declare where in the
term those rewrites happen and in what order.
\begin{Verbatim}
reify-eval (c-fmap {a} {b}) =
   c-trans (c-sym c-arr-fmap)
    (c-trans
      (c-cong-comp
        (c-refl c-fmap)
        (c-trans
          (c-cong-arr (c-trans (c-cong-mono (reify-mkid a)) c-mono-id)
                      (c-trans (c-cong-mono (reify-mkid b)) c-mono-id))
          c-arr-id))
      c-id-r)
\end{Verbatim}



\section{Perspectives}
Appealing extensions: applicatives, monads, \dots

Difficulties in doing that: loss of the peel invariant, etc

Appealing results: coherence for an entire lambda calculus

Difficulties in doing that: difficult to set up an untyped semantics in order to
prove a bracketing theorem

\section{Related Work}
Reynolds~\cite{reynolds80}.

Coherence proof for monoidal categories~\cite{beylin95}.

\section{Conclusion}

\bibliographystyle{ACM-Reference-Format}
\bibliography{paper.bib}

\end{document}
\endinput
%%
%% End of file `sample-acmsmall-submission.tex'.
